\head{第三版新核查：第三卷第二部分及作者勘误}
\noindent\textbf{结论先行。}本轮不是对全部已译文字重新作一遍数学审定，而是在续译第三卷（6.7.1）--（7.9.14）及其书末勘误时作定点复核。可以确定：原（7.3.10）中不加分支限制的“在某个稠密开集上恒定”有反例；（6.7.10）、（7.4.2）、（7.4.6）的若干中间说法也不成立，但这不等于相应主结论全部错误。作者自己已经刊出的勘误另列，不当作本次原创发现。

本章的反例、补证和处理建议均为译校分析。被翻译的原文仍保留其组织与证明详略，不能把本章补证归入原作者正文。以下“原页”指原分册印刷页码；上传PDF页码比正文原分册页码少1。末尾勘误另有连续总页码217--223，均在定位处写明。

\section*{定位与状态表}
\addcontentsline{toc}{section}{第三版新核查定位表}
\begin{longtable}{@{}p{.17\linewidth}p{.24\linewidth}p{.51\linewidth}@{}}
\toprule 标识 & 原条目／PDF页 & 本轮判定\\\midrule\endhead
C3-E01 & III 7.3.10／51 & 全局“恒定”断言在所写假设下有反例；不否定局部恒定等较弱结论。\\
C3-G01 & III 6.7.10／30 & “子双复形”的注入性无保证；典范同态可以有核。\\
C3-G02、G03 & III 7.3.10／51 & 底环与商环的归约需区分；局部环等同错误。\\
C3-G04 & III 7.4.2／54 & 一条短正合列的末项错置，并泛称循环与张量交换；主结论有修补证明。\\
C3-G05 & III 7.4.6／55 & 核的对偶不一般等于转置的余核；表示函子结论有直接补证。\\
C3-P01 & III 6.7.1--2／25 & 解消的投射性所属底环不明确；不能混用投射性。局部化有替代论证。\\
C3-H01 & III 6.7.3／26 & 仿射计算式须继承先在共同底仿射开片计算的语境。\\
C3-T01--T09 & 见相应脚注 & 覆盖对象、引文、次数与符号问题；不统一计作错误定理。\\
C3-A01--A03 & 原作者勘误／86--88 & 正式修订与旧稿状态更新；另说明两处补证细节。\\
\bottomrule
\end{longtable}

\audit{C3-E01}{稠密开集上的恒定不等于局部恒定}{原表述有明确反例}
\source{EGAIII.2.pdf，第51页，原第52页，命题（7.3.10）}
该条假设 $A$ 诺特，同调函子 $T_\bullet$ 与归纳极限交换，各 $T_i(A)$ 有限型，并在某个次数以下正合。它声称，对给定上界 $p$，存在稠密开集 $E\subset\Spec A$，使所有 $i\leq p$ 的函数 $x\mapsto\dim_{k(x)}T_i(k(x))$ 在 $E$ 上恒定。

取任意域 $k$，令
\[
 A=k\times k,\qquad e=(1,0),\qquad
 T_0(M)=eM,\qquad T_i(M)=0\ (i\ne0).
\]
$A$ 是诺特环；$eA$ 是 $A$ 的直和因子，故 $M\mapsto eM$ 正合并与归纳极限交换。将所有连接同态取为零，$T_\bullet$ 是同调函子：任意短正合列在第0项仍给出短正合列，其他次数均为零。$T_0(A)=eA$ 由 $e$ 生成，其他 $T_i(A)=0$，满足所要求的有限型条件。事实上，各 $T_i$ 在所有次数都正合，因而也满足“充分低次数正合”的假设。

空间 $\Spec A$ 有两个开闭点。对应商映射 $A\to k$ 的第一坐标点记作 $x_1$，第二坐标点记作 $x_2$，则
\[
 T_0(k(x_1))=k,
 \qquad T_0(k(x_2))=0.
\]
所以秩函数分别取1和0。有限离散空间的稠密开集只能是全空间，因此不存在使此函数恒定的稠密开集。

\textbf{修订范围。}应当区分“每个不可约分支中的适当稠密开集上恒定”和“整个稠密开集上有同一个值”。原证明在各分支分别作选择，不能迫使不同分支上的常数相等。本例的秩函数在全空间已经局部恒定，却仍不恒定。若使用弱化为局部恒定的版本，还须补足下一条所述归约步骤，不能仅换一个词就视为原证明每一步均已正确。

\audit{C3-G02--G03}{商环归约与局部环不能混同}{证明中的归约缺步及中间等式错误}
\source{EGAIII.2.pdf，第51页，原第52页，（7.3.10）证明}
原证明先取 $B=A/\mathfrak a$ 为不可约分支的坐标环，在 $\Spec B$ 的适当开集证明相应函子性质，再写出将 $A$ 局部环与 $B$ 局部环等同的步骤。对商映射 $A\to B$，点可以对应，剩余域也可以对应，但一般有
\[
 B_{\mathfrak p/\mathfrak a}
 =A_{\mathfrak p}/\mathfrak a A_{\mathfrak p},
\]
而不是 $B_{\mathfrak p/\mathfrak a}=A_{\mathfrak p}$。

例如 $A=k[\epsilon]/(\epsilon^2)$、$B=A/(\epsilon)=k$。两者的谱都只有一个点，且都是整个空间的稠密开集，但 $A$ 有非零幂零元而 $B$ 没有；取 $f=1$ 后，局部化也不能使它们成为同一个环。若原论证只想比较剩余域上的函数值，应明确只使用 $k(x)$ 的一致性，不可借此搬运整个局部环上的正合性。

此外，在 $B$ 上得出一个次数的正合性，并不能未经说明便将 $A$ 上的“已知正合次数下界”提高到该次数。仍取 $A=k[\epsilon]/(\epsilon^2)$、$B=k$，并取同调函子 $T_i(M)=\Tor_i^A(k,M)$，负次数取零。它在负次数正合，与归纳极限交换，且 $T_0(A)=k$、$T_i(A)=0$（$i>0$）均有限型。在 $B$-模范畴，$T_0^{(B)}(M)=k\otimes_A M=M$ 正合；但在 $A$-模范畴，$T_0$ 不左正合：单射 $(\epsilon)\hookrightarrow A$ 在张量 $k$ 后成为 $k\to k$ 的零映射。这具体显示原文在归纳过程中提升正合范围的步骤不能仅凭商环上的结论完成。该部分须另作逐分支归纳论证；本轮不据此否定更弱的可构造性或局部恒定结论。

\audit{C3-G01}{典范Čech映射不一定是子复形包含}{中间说法错误，主构造仍可用同态表述}
\source{EGAIII.2.pdf，第30页，原第31页，（6.7.10）}
原文在比较基变换的Čech双复形时使用“子双复形”一词。这要求逐项单射，而一般的拉回截面同态没有这一性质。

令 $S=Y=Y'=\Spec k$，$X=\Spec k[t]$，$X'=\Spec k$，$u:X'\to X$ 由 $t\mapsto0$ 给出；取 $\F=\OO_X$，各空间都用单个仿射开集覆盖。于是涉及的第0项典范映射是
\[
 k[t]=\Gamma(X,\OO_X)
 \longrightarrow\Gamma(X',u^*\OO_X)=k,
 \qquad p(t)\longmapsto p(0).
\]
其核包含且等于非零理想 $(t)$。故一般只能说有典范双复形同态，不能称其为子双复形包含。随后若仅用该同态诱导超同调或谱序列映射，仍有可能完成构造；这个反例不否定映射的存在，也不否定在额外平坦条件下的同构断言。

\audit{C3-G04}{（7.4.2）中短正合列及张量交换的修补}{证明步骤错误；给出不改变主结论的补证}
\source{EGAIII.2.pdf，第54页，原第55页，（7.4.2）证明}
记 $d_i:P_i\to P_{i-1}$，$Z_i=\Ker d_i$，$B_i=\Im d_{i+1}$，$Z'_i=P_i/B_i$。在需要构造替代平坦复形的步骤，原文使用
\[
 0\to Z_{p+1}\to P_{p+1}\to Z'_p\to0.
\]
但 $P_{p+1}\xrightarrow{d_{p+1}}P_p$ 的像是 $B_p$，并非 $Z'_p$；从定义得到的是
\[
 0\to Z_{p+1}\to P_{p+1}\to B_p\to0.
\]

\textbf{同时反驳原列及泛称张量交换。}取 $A=k[t]$，复形仅有 $P_1=A\xrightarrow{t}P_0=A$，并取 $p=1$。所有项平坦，且 $Z'_1=A$ 平坦，符合该步骤的假设。然而原列成为 $0\to0\to0\to A\to0$，不正合。又
\[
 Z_1(P_\bullet)=0,
 \qquad Z_1(P_\bullet\otimes_A k)=k
 \quad(k=A/(t)),
\]
所以即使各项平坦，循环模也不一般与任意张量交换。

\textbf{所需的正确局部论证。}假设 $P_i$ 均平坦，且 $Z'_p$ 平坦。由
\[
 0\to B_p\to P_p\to Z'_p\to0
\]
可知 $B_p$ 平坦：对任意 $A$-模 $N$，长Tor列中的 $\Tor_2(Z'_p,N)$ 与 $\Tor_1(P_p,N)$ 均为零，故 $\Tor_1(B_p,N)=0$。再由
\[
 0\to Z_{p+1}\to P_{p+1}\to B_p\to0
\]
同理得到 $Z_{p+1}$ 平坦。两条列的余核都平坦，故它们在张量任意 $M$ 后仍正合。

定义新复形 $P'_\bullet$：在次数 $p$ 取 $Z'_p$，次数 $p+1$ 取 $Z_{p+1}$，其余次数不变。$d'_{p+1}=0$；$d'_p:Z'_p\to P_{p-1}$ 是 $d_p$ 降到商模后的同态；$d'_{p+2}$ 是 $d_{p+2}$ 因其像落在 $Z_{p+1}$ 而得到的余限制。其余微分不变。刚证明的平坦性说明新复形各项仍平坦。

对任意 $M$，在次数 $p$，商映射 $P_p\otimes M\to Z'_p\otimes M$ 给出
\[
 \frac{\Ker(d_p\otimes1)}{\Im(d_{p+1}\otimes1)}
 \simeq\Ker(d'_p\otimes1).
\]
原因是商映射满射，其核恰为 $B_p\otimes M$，而 $P_{p+1}\otimes M\to B_p\otimes M$ 满射。在次数 $p+1$，两条张量后的正合列又给出
\[
 \Ker(d_{p+1}\otimes1)=Z_{p+1}\otimes M,
\]
故新旧同调都等于此模除以 $d_{p+2}\otimes1$ 的像。次数 $p-1$ 的边界像不变；次数 $p+2$ 的循环不变，因为 $Z_{p+1}\otimes M\to P_{p+1}\otimes M$ 仍单射。其余次数的复形项与微分均不变。因此所有同调模都有关于 $M$ 的自然同构。

连接同态也相容：对 $0\to M'\to M\to M''\to0$，连接同态按“提升循环、施加微分、取回左端”定义；上述同构在次数 $p$ 由商映射给出，在次数 $p+1$ 由保留单射的包含给出，二者与微分交换，所以同一提升得到相同的左端同调类。两者之间的连接同态为零，亦与 $d'_{p+1}=0$ 一致。这个论证只在所需的两个次数使用张量正合性，没有声称所有 $Z_i,B_i$ 都与张量交换。

\audit{C3-G05}{核的对偶与转置余核不同}{原证明的辨认错误；表示函子结论保留}
\source{EGAIII.2.pdf，第55页，原第56页，（7.4.6）}
对同态 $d:P\to R$，一般不能把 $(\Ker d)^\vee$ 认作 $\Coker(d^\vee)$。仍取 $A=k[t]$、$P=R=A$、$d$ 为乘 $t$。此时
\[
 (\Ker d)^\vee=0,
 \qquad\Coker(d^\vee)=A/(t)\ne0.
\]
所以这两个模不同。

\textbf{正确的表示模。}应用上一条对（7.4.2）的修补后，所需函子可写成 $M\mapsto\Ker(P\otimes_A M\to R\otimes_A M)$，其中 $P,R$ 为有限型投射模。直接定义
\[
 Q=\Coker(R^\vee\xrightarrow{d^\vee}P^\vee).
\]
对正合列 $R^\vee\to P^\vee\to Q\to0$ 应用左正合的反变函子 $\Hom_A(-,M)$，得到
\[
 0\to\Hom_A(Q,M)\to\Hom_A(P^\vee,M)
 \to\Hom_A(R^\vee,M).
\]
有限型投射模的典范同构 $\Hom_A(P^\vee,M)\simeq P\otimes_A M$ 及 $\Hom_A(R^\vee,M)\simeq R\otimes_A M$ 将最右同态识别为 $d\otimes1$；因此
\[
 \Hom_A(Q,M)\simeq\Ker(d\otimes1)
\]
且关于 $M$ 自然。$Q$ 有限呈示，诺特条件下当然有限型。例子中 $Q=A/(t)$，且 $\Hom_A(A/(t),M)=\{m\in M:tm=0\}$，正好是所需核函子。整个证明无需声称 $Q=(\Ker d)^\vee$。

表示对象的唯一性应理解为连同指定的表示同构，唯一到与该同构相容的唯一同构；并不表示裸模没有非平凡自同构。

\audit{C3-P01}{投射解消的底环与局部化}{原步骤须补明；局部化断言不因此被否定}
\source{EGAIII.2.pdf，第25页，原第26页，（6.7.1）、（6.7.2）}
原文先计算 $A$ 上的超Tor，继而选择“由 $B_i$-模组成的投射 Cartan--Eilenberg 解消”。若“投射”指 $A$-投射，一般不存在这种选择。例如 $A=\mathbf Z$、$B_i=\mathbf F_p$ 时，任何非零 $B_i$-模都被 $p$ 零化，不可能是 $\mathbf Z$-投射模；若改为 $B_i$-投射，则又不能直接用它计算 $A$-导出张量积。

后一差别可以直接计算：$\mathbf F_p$ 在自身上自由，而
\[
 \Tor_1^{\mathbf Z}(\mathbf F_p,\mathbf F_p)=\mathbf F_p.
\]
取 $\mathbf F_p$ 在自身上的零长度解消，会把这个群错误地算成零。因此必须注明解消范畴和相对非零Tor，而不能只靠术语“投射”。

\textbf{避开该选择的一条路径。}在（6.7.1）的仿射计算中，对双复形 $L^{(i)}$，局部化可写为
\[
 L^{(i)}_{g_i}=\varinjlim\bigl(L^{(i)}\xrightarrow{g_i}L^{(i)}
 \xrightarrow{g_i}\cdots\bigr).
\]
双索引系统 $\mathbf N\times\mathbf N$ 的对角子系统共尾。应用（6.3.6）中超Tor与滤过归纳极限的交换性，得到
\[
 \operatorname{Tor}_n^A(L^{(1)}_{g_1},L^{(2)}_{g_2})
 \simeq\varinjlim_m\operatorname{Tor}_n^A(L^{(1)},L^{(2)}),
\]
右侧过渡映射由两变量分别乘 $g_1,g_2$ 诱导，因双函子性正是乘 $g_1\otimes g_2$。所以右侧是原超Tor的 $(g_1\otimes g_2)$-局部化。

这里的Čech方向只有有限多非零次数，另一方向有统一下界，故总复形与上式所需的有限次数操作相容。六个谱序列的 $E^2$ 表达式由同调、直和及Tor构成；滤过归纳极限正合，且这些操作依（6.3.6）相容，于是同样得到相应 $E^2$ 页的比较同构。共同终值的同构已如上得到；在原文的双正则条件下，可作相应谱函子的比较。该说明给出替代局部化论证的必要步骤，不把原文不明确的投射解消选择当作已经合法。

（6.7.2）中从 $A$ 改为 $A_h$，应使用平坦局部化后的超Tor比较（6.3.13），而非声称每个 $A_h$-投射模都是 $A$-投射。原文已有的局部化引理提供比较依据。

\audit{C3-H01}{一般底空间上的仿射图}{继承局部化语境，不作无条件计算}
\source{EGAIII.2.pdf，第26页，原第27页，（6.7.3）}
当 $S$ 任意时，不能仅凭 $V^{(1)},V^{(2)}$ 仿射，就无条件把 $V^{(1)}\times_SV^{(2)}$ 当作仿射概形，再对一个环模加波浪号。本条前文的构造先在 $S$ 上取共同仿射开集，再取两侧其上方的仿射开图，正是在这个语境下才能使用“模的相伴层”的写法。译文保留原式，脚注明确它应按前述局部计算与粘合解释。

\section*{次数、符号及引文记录}
\addcontentsline{toc}{section}{第三版次数、符号及引文记录}
\begin{longtable}{@{}p{.17\linewidth}p{.24\linewidth}p{.51\linewidth}@{}}
\toprule 标识 & 原位置 & 处理与理由\\\midrule\endhead
C3-T01 & 6.7.4，PDF26 & 用有限仿射覆盖计算源空间上同调时，应覆盖 $X$；该句写为覆盖 $Y$。保留正文并注。\\
C3-T02 & 6.9.5 & 准同构或同伦等价不变性的直接依据为（6.7.5）；原文引用（6.7.6），后者是平坦复形比较。\\
C3-T03 & 6.10.5 & 选取有限自由项所用的是（6.10.3，iii）；比较同构部分才用（ii）。\\
C3-T05 & 7.9.2，PDF75 & 证明中（7.9.1.2）应指紧邻的（7.9.2.1）。\\
C3-T06 & 7.9.4，PDF75 & 按（7.7.1.1），同调指标 $p$ 对应超上同调指标 $-p$。译稿显式标明统一次数约定，不抹去原印刷差异。\\
C3-T07 & 7.9.6，PDF76 & 图中复形箭头的方向与同调次数递减约定不一致；图依底本保留，截断证明按次数约定解释。\\
C3-T08 & 7.9.7，PDF77 & 引用（6.10.4.2）未见对应编号；基变换的函数值相等由（7.7.2）与（7.6.7）得到。\\
C3-T09 & 7.9.8，PDF77 & 陈述中的秩为 $(-1)^{i_0}\operatorname{EP}$ 正确；证明末行漏掉这个符号。\\
\bottomrule
\end{longtable}
上述编号与译文脚注对应，编号并不保证连续，也不是错误定理的计数器。C3-N开头的脚注说明同调次数、左右模、半正合和逐项Tor等记号；C3-H02记录可通过放大界限吸收的端点差异；C3-H03将原著提出的问题保留为历史问题，不凭本轮翻译宣布其今天的研究状态。

\textbf{C3-T09的直接检验。}若只在 $i_0$ 次同调非零，则定义给出
\[
 \operatorname{EP}=\sum_i(-1)^i\dim H_i=(-1)^{i_0}\dim H_{i_0},
\]
故 $\dim H_{i_0}=(-1)^{i_0}\operatorname{EP}$。例如仅在次数1有一个一维向量空间时，秩为1而Euler特征数为 $-1$；遗漏符号的等式就错误。原命题陈述已经含正确符号，不需推翻命题。

\audit{C3-A01}{把正式勘误接回正文}{原作者修订，不是新发现}
\source{EGAIII.2.pdf，第84--90页，《勘误与增补（清单二）》}
本轮完整译入作者勘误的A、B两部分，B有30个编号条目。下面列的是与前次审阅紧密相关的几项，不是对这30条的摘要替代；全文已在合订本书末。

\begin{longtable}{@{}p{.19\linewidth}p{.24\linewidth}p{.49\linewidth}@{}}
\toprule 正式勘误 & 所修正位置 & 原作者所作处理\\\midrule\endhead
ErrIII，3 & $0_{\mathrm I}$ 7.2.4--7.2.6 & 增加各理想幂为闭理想的假设；指出原证明不正确。\\
ErrIII，10 & I 5.3.9 & 原证明不足，补充仿射局部的对角闭浸入论证。\\
ErrIII，19 & II 5.2.2 & 局部性论证不充分，改用 $X$ 是概形或底空间局部诺特的假设。\\
ErrIII，21 & II 6.4.3 & 否定“有限型子模元素必整”的推理，改用忠实有限模证明自同态整。\\
ErrIII，24 & $0_{\mathrm{III}}$ 13.7.7 & 改正ML与有限性结论以及相关次数范围。\\
\bottomrule
\end{longtable}
前几版某些译注只指出疑点或给出独立反例，未附这份正式勘误。本版在对应条目前添加跳转提示。旧审阅报告作为版本历史保留，但读者应以这里给出的作者勘误状态为准，不能将已有正式修订的问题继续宣传为“尚无人修正”或本次的新发现。

\textbf{闭理想条件为何足够。}设 $A$ 对原拓扑分离完备，$J^n$ 趋于0且各 $J^n$ 在原拓扑闭。给定 $\varprojlim A/J^n$ 的元素，选相容类的代表 $a_n$，使 $a_m-a_n\in J^n$（$m\geq n$）。因 $J^n\to0$，序列 $(a_n)$ 对原拓扑是柯西列，故收敛到某个 $a\in A$。固定 $n$，尾项 $a_m-a_n$ 全在闭集 $J^n$，于是其极限 $a-a_n\in J^n$。所以 $a$ 表示全部给定相容类。分离性及 $J^n\to0$ 又给出 $\bigcap_nJ^n=0$，保证唯一性。故 $A\to\varprojlim A/J^n$ 双射，这证明了修订条件下的 $J$-进制分离完备性。闭性正用于从“尾项属于理想”推出“极限属于理想”。

\audit{C3-A02}{对角态射补证中的共同仿射底}{补足原勘误的局部化说明}
\source{EGAIII.2.pdf，第87页，ErrIII，10；EGA I（5.3.9）}
对 $x\in X$，先取 $f(x)$ 在 $S$ 中的仿射开邻域 $V=\Spec B$，再取仿射开邻域 $x\in U=\Spec A\subset f^{-1}(V)$。此时
\[
 W_U=U\times_VU=U\times_SU
\]
是 $X\times_SX$ 的开子概形，而且仿射；$\Delta^{-1}(W_U)=U$。受限对角态射对应于乘法满同态
\[
 A\otimes_BA\longrightarrow A,
 \qquad a\otimes a'\longmapsto aa'.
\]
它满射，因为 $a$ 是 $a\otimes1$ 的像。因此 $U\to W_U$ 是闭浸入。

令 $W$ 为所有这些 $W_U$ 的并，则 $W$ 是包含整个对角像的开子概形。在每个 $W_U$ 上，$\Delta(X)\cap W_U=\Delta(U)$ 闭，故对角像在 $W$ 中闭；相应核理想在重叠处相容，因为都来自同一结构层同态。因此这些局部闭浸入粘合为 $X\to W$ 的闭浸入。再复合 $W\hookrightarrow X\times_SX$ 的开浸入，得到所需浸入。

这说明原作者指出的缺步可以补足；也说明不能只写“取任意仿射 $U$，所以 $U\times_SU$ 仿射”而漏掉先选共同仿射底的步骤。

\audit{C3-A03}{可逆层证明中的非零芽与非零值}{文字条件须加强，结论可补证}
\source{EGAIII.2.pdf，第86页，ErrIII，2的续文}
作者补证中写“可设 $s=s_1$ 且 $s_x\ne0$”，随后由 $\LL_x\simeq\OO_x$ 推出它生成该模。单凭非零芽不足：在 $A=k[t]_{(t)}$ 上，$t\ne0$，却不生成秩一自由模 $A$。

这里已有有限生成元 $s_1,\ldots,s_n$。将它们映入一维剩余向量空间 $\LL_x/\mathfrak m_x\LL_x$，至少一个像非零，否则它们不能生成 $\LL_x$。重新编号，选这个元素为 $s_1$；其在局部自由秩一模中的系数是单位，故确实生成 $\LL_x$。于是后面的局部截面 $t_i$ 可按原证明构造。应写明“值 $s(x)\ne0$”或“$s_x$ 生成”，而不是把任意非零芽视为生成元。译文保留原词并作脚注，不静默替换。

\section*{对本版译文的核对范围}
\addcontentsline{toc}{section}{第三版核对范围与限制}
本轮逐段续译第三卷第二部分剩余正文，并核对末页实际收录到（7.9.14）及其证明、索引、目录、勘误。新增数学反例只覆盖本章列出的断言，不能推及未检查的全部证明。原作省略的证明、历史问题以及英语预备笔记的未完成内容，仍明确标示。

本次重排保留旧正文的条目与公式号，把第0章补篇归并、第三章上下部分连续排入，将书末资料移至正文之后。该结构操作与编译成功均不是数学正确性的证明。第四卷主体、第0章§§16--23和第V卷余稿尚未连续译入；其原文的存在、原目录的译名、文件转换成功，都不算翻译完成。

本报告后半部保存先前核查文字作为历史档案。本章没有重新验证其中每一个反例，也没有重新核查其全部外部文献。遇到与本章新取得的原作者勘误状态不同的旧表述，应同时阅读本章更新说明。
